\section{Séparateurs (Separator)}

\subsection{Structure \texttt{separator\_config}}
\begin{lstlisting}[language=C]
typedef struct {
    GtkOrientation orientation;
    widget_style_config style;
} separator_config;
\end{lstlisting}

\subsection{Description des champs}
\begin{center}
\begin{tabular}{|l|l|>{\raggedright\arraybackslash}p{6cm}|}
\hline
\textbf{Champ} & \textbf{Type} & \textbf{Description} \\
\hline
orientation & GtkOrientation & \texttt{GTK\_ORIENTATION\_HORIZONTAL} ou \texttt{VERTICAL} \\
\hline
style & widget\_style\_config & Configuration (marges, CSS) \\
\hline
\end{tabular}
\end{center}

\bigskip

Le champ \texttt{orientation} permet de définir si la ligne de séparation doit être horizontale ou verticale (\texttt{GTK\_ORIENTATION\_HORIZONTAL} ou \texttt{VERTICAL}). C'est l'outil de design par excellence pour délimiter visuellement deux zones distinctes dans une interface chargée.

On l'utilise couramment pour séparer une barre d'outils du contenu principal ou pour diviser deux colonnes de texte, aidant ainsi l'utilisateur à comprendre la structure hiérarchique de l'application.

Le champ \texttt{style} permet d'ajuster l'apparence du séparateur via la structure \texttt{widget\_style\_config}. Il est particulièrement utile pour ajouter des marges (\texttt{margin}) afin que la ligne ne touche pas les bords des widgets adjacents, créant une interface plus aérée et professionnelle.